\documentclass[11pt]{article}
\usepackage{mathblatt}
% Im Vorspann ergänzt (fehlt in der Anleitung): leeres Gitter mit Achsenpfeilen, ohne Bezifferung,
% Karo 8 mm – für Aufgaben, in denen der Schüler die Achseneinteilung selbst wählt (Nr. 25).
\newcommand{\leeresgitter}[2]{\par\noindent\begin{tikzpicture}
\draw[step=0.8cm,gray!45,very thin] (0,0) grid ({#1*0.8},{#2*0.8});
\draw[->,thick] (0,0) -- ({#1*0.8+0.5},0);
\draw[->,thick] (0,0) -- (0,{#2*0.8+0.5});
\end{tikzpicture}\par}
\newif\ifmitzone
\mitzonetrue
\begin{document}
\blattkopf{Potenz- und Exponentialfunktionen}{Gesamt}
\weit
\verzeichniszeile{\verz{zone}{Kennst du schon (Nr.~1–12)} \verztrenn \verz{e1}{1 Linear oder exponentiell (Nr.~13–29)} \verztrenn \verz{e2}{2 Wachstumsfaktor und Tabelle (Nr.~30–39)} \verztrenn \verz{e3}{3 Exponentialfunktion (Nr.~40–49)} \verztrenn \verz{e4}{4 Verdopplungs- und Halbwertszeit (Nr.~50–57)} \verztrenn \verz{e5}{5 Potenzfunktionen (Nr.~58–69)} \verztrenn \verz{abhaken}{Das kann ich}}

% Zone „Kennst du schon" – Nr. 1–12
\setcounter{aufgabe}{0}
\einheitenkopf[zone][Kennst du schon]{Das kennst du schon}

\begin{aufgabe}{Ich kann Dezimalzahlen multiplizieren und das Ergebnis runden.}
Rechne.
\begin{teilezwei}
\tz $4 \cdot 1{,}2$ \leerfeld & \tz $30 \cdot 0{,}9$ \leerfeld \\
\tz $250 \cdot 1{,}04$ \leerfeld & \tz $0{,}8 \cdot 0{,}8$ \leerfeld \\
\tz $7{,}5 \cdot 1{,}12$ \leerfeld & \\
\anweisung{Runde auf eine Nachkommastelle.}
\tz $45 \cdot 1{,}15 \cdot 1{,}15$ \leerfeld & \tz $1{,}8 \cdot 2{,}35$ \leerfeld \\
\end{teilezwei}
\end{aufgabe}

\begin{aufgabe}{Ich kann Punkte in ein Koordinatensystem eintragen.}
Trage die Punkte ein und beschrifte sie.
\begin{teile}
\teil $A(1\,|\,2)$
\teil $B(4\,|\,5)$
\teil $C(2{,}5\,|\,3{,}5)$
\end{teile}

\begin{ksys}[xmin=0,xmax=6,ymin=0,ymax=6]
\end{ksys}

\end{aufgabe}

\begin{aufgabe}{Ich kann Punkte eintragen, wenn eine Achse in größeren Schritten eingeteilt ist.}
Trage die Punkte ein und beschrifte sie.
\begin{teile}
\teil $D(1\,|\,100)$
\teil $E(3\,|\,250)$
\teil $F(4\,|\,375)$
\end{teile}

\begin{ksys}[xmin=0,xmax=5,ymin=0,ymax=400,ystep=50]
\end{ksys}

\end{aufgabe}

\begin{aufgabe}{Ich kann eine passende Einteilung für eine Achse wählen.}
Die senkrechte Achse hat eine bestimmte Zahl von Kästchen. Kreuze an, welche Einteilung je Kästchen die Achse am besten ausnutzt.
\begin{teile}
\teil Werte bis 40, 8 Kästchen: \quad \kreuz{1}\kreuz{5}\kreuz{10}
\teil Werte bis 900, 10 Kästchen: \quad \kreuz{50}\kreuz{100}\kreuz{200}
\teil Werte bis 2000, 10 Kästchen: \quad \kreuz{100}\kreuz{200}\kreuz{500}
\teil Werte bis 3,5, 7 Kästchen: \quad \kreuz{0,5}\kreuz{1}\kreuz{2}
\end{teile}
\end{aufgabe}

\begin{aufgabe}{Ich kann einen Wert um einen Prozentsatz erhöhen oder senken.}
Berechne den neuen Wert.
\begin{teile}
\teil 300\,€ werden um 10\,\% erhöht. \leerfeld[€]
\teil 50\,kg werden um 20\,\% gesenkt. \leerfeld[kg]
\teil 64\,€ werden um 5\,\% erhöht. \leerfeld[€]
\teil 180\,€ werden um 2,5\,\% gesenkt. \leerfeld[€]
\teil 45\,€ werden auf 80\,\% gesenkt. \leerfeld[€]
\teil 200\,€ werden erst um 10\,\% und danach noch einmal um 10\,\% erhöht. \leerfeld[€]
\end{teile}
\end{aufgabe}

\begin{aufgabe}{Ich kann eine Zahlenfolge fortsetzen, wenn immer gleich viel dazukommt.}
Setze die Folge um zwei Zahlen fort.
\begin{teile}
\teil 5, \ 8, \ 11, \ \leerfeld, \ \leerfeld
\teil 40, \ 35, \ 30, \ \leerfeld, \ \leerfeld
\teil 1,2; \ 1,5; \ 1,8; \ \leerfeld; \ \leerfeld
\end{teile}
\anweisung{Gib die Werte an.}
\begin{teile}
\teil Ein Taxi kostet 4\,€ Grundpreis und 2\,€ je Kilometer. Preis für 0\,km, 1\,km und 5\,km: \\ \leerfeld[€] \leerfeld[€] \leerfeld[€]
\teil An welcher Stelle trifft die Gerade $y = 3x + 5$ die senkrechte Achse? \leerfeld
\end{teile}
\end{aufgabe}

\begin{aufgabe}{Ich kann zwei Zahlen durcheinander teilen und das Ergebnis als Dezimalzahl angeben.}
Berechne den Quotienten.
\begin{teilezwei}
\tz $12 : 4$ \leerfeld & \tz $50 : 25$ \leerfeld \\
\tz $45 : 30$ \leerfeld & \tz $54 : 60$ \leerfeld \\
\anweisung{Teile den zweiten Wert durch den ersten.}
\tz von 80 auf 100: \leerfeld & \tz von 250 auf 200: \leerfeld \\
\end{teilezwei}
\end{aufgabe}

\begin{aufgabe}{Ich kann Prozentwert und Prozentsatz berechnen.}
Berechne den Prozentwert.
\begin{teilezwei}
\tz 10\,\% von 70\,€ \leerfeld[€] & \tz 25\,\% von 80\,kg \leerfeld[kg] \\
\tz 3\,\% von 450\,€ \leerfeld[€] & \\
\anweisung{Gib an, wie viel Prozent das sind.}
\end{teilezwei}
\begin{teile}
\teil 18 von 120 Schülern \leerfeld[\%]
\teil Ein Preis steigt von 50\,€ auf 56\,€. Um wie viel Prozent ist er gestiegen? \leerfeld[\%]
\end{teile}
\end{aufgabe}

\begin{aufgabe}{Ich kann ein Guthaben mit Zinseszins berechnen.}
Die Zinsen werden jedes Jahr zum Guthaben addiert. Berechne das Guthaben.
\begin{teile}
\teil 1000\,€ zu 2\,\% Zinsen, nach einem Jahr: \leerfeld[€]
\teil dasselbe Guthaben nach zwei Jahren: \leerfeld[€]
\teil 500\,€ zu 4\,\% Zinsen, nach drei Jahren: \leerfeld[€]
\end{teile}
\end{aufgabe}

\begin{aufgabe}{Ich kann Potenzen mit dem Taschenrechner berechnen.}
Berechne.
\begin{teilezwei}
\tz $2^5$ \leerfeld & \tz $10^3$ \leerfeld \\
\tz $1{,}5^3$ \leerfeld & \tz $(-2)^3$ \leerfeld \\
\anweisung{Runde auf zwei Nachkommastellen.}
\tz $0{,}9^4$ \leerfeld & \tz $3 \cdot 1{,}2^4$ \leerfeld \\
\end{teilezwei}
\end{aufgabe}

\begin{aufgabe}{Ich finde den Fehler in einer Zinseszinsrechnung.}
Tom legt 800\,€ für drei Jahre zu 5\,\% Zinsen an. Die Zinsen werden jedes Jahr zum Guthaben addiert. Tom rechnet so:
\rechnung{3 \cdot 5\,\% &= 15\,\% \\ 800\,€ \cdot 1{,}15 &= 920\,€}
Finde den Fehler, benenne ihn und korrigiere die Rechnung.
\end{aufgabe}

\begin{aufgabe}{Ich kann das Guthaben nach mehreren Jahren Zinseszins berechnen.}
Berechne das Guthaben.
\begin{teile}
\teil 1500\,€ zu 4\,\% Zinsen, nach zwei Jahren: \leerfeld[€]
\end{teile}
\end{aufgabe}

\clearpage
% Einheit 1 – Nr. 13–29
\setcounter{aufgabe}{12}
\einheitenkopf[e1][1 Linear oder exponentiell]{Einheit 1 von 5 · Lineares und exponentielles Wachstum unterscheiden}
\zweigzeile{Hier lernst du, lineares und exponentielles Wachstum an Tabelle, Text und Graph zu unterscheiden · neu in diesem Jahr · P10 oft · baut auf: Punkte eintragen, gleicher Betrag je Schritt, Erhöhen um p\,\%}

\begin{aufgabe}{Ich erkenne, ob jedes Mal derselbe Betrag oder derselbe Prozentsatz dazukommt.}
Kreuze an. Rechne nichts.
\begin{teile}
\teil Ein Sparschwein bekommt jeden Monat 15\,€ dazu. \\ \kreuz{gleicher Betrag}\kreuz{gleicher Prozentsatz}
\teil Ein Guthaben wächst jedes Jahr um 2\,\%. \\ \kreuz{gleicher Betrag}\kreuz{gleicher Prozentsatz}
\teil Eine Bohnenpflanze wird jeden Tag 3\,cm höher. \\ \kreuz{gleicher Betrag}\kreuz{gleicher Prozentsatz}
\teil Die Zahl der Nutzer einer App steigt jeden Monat um 12\,\%. \\ \kreuz{gleicher Betrag}\kreuz{gleicher Prozentsatz}
\teil Ein Auto verliert jedes Jahr 15\,\% seines Werts. \\ \kreuz{gleicher Betrag}\kreuz{gleicher Prozentsatz}
\end{teile}
\end{aufgabe}

\begin{aufgabe}{Ich erkenne, ob in einer Tabelle die Differenz oder der Quotient gleich bleibt.}
Unter den Werten stehen schon die Differenzen und die Quotienten benachbarter Werte. Kreuze an, was gleich bleibt.
\begin{teile}
\teil Werte: 4; \ 12; \ 36; \ 108 \quad Differenzen: 8; \ 24; \ 72 \quad Quotienten: 3; \ 3; \ 3 \\ \kreuz{die Differenz}\kreuz{der Quotient}
\teil Werte: 50; \ 65; \ 80; \ 95 \quad Differenzen: 15; \ 15; \ 15 \quad Quotienten: 1,3; \ $\approx 1{,}23$; \ $\approx 1{,}19$ \\ \kreuz{die Differenz}\kreuz{der Quotient}
\teil Werte: 400; \ 200; \ 100; \ 50 \quad Differenzen: $-200$; \ $-100$; \ $-50$ \quad Quotienten: 0,5; \ 0,5; \ 0,5 \\ \kreuz{die Differenz}\kreuz{der Quotient}
\teil Werte: 90; \ 81; \ 72; \ 63 \quad Differenzen: $-9$; \ $-9$; \ $-9$ \quad Quotienten: 0,9; \ $\approx 0{,}89$; \ $\approx 0{,}88$ \\ \kreuz{die Differenz}\kreuz{der Quotient}
\end{teile}
\end{aufgabe}

\begin{aufgabe}{Ich erkenne, wo ein Graph die senkrechte Achse trifft.}
Markiere den Punkt, in dem der Graph die senkrechte Achse trifft. Rechne nichts.

\begin{ksys}[xmin=0,xmax=5,ymin=0,ymax=8,ystep=2,ablesen]
\gerade{1.2}{0}{A}
\end{ksys}\ksysabstand\begin{ksys}[xmin=0,xmax=5,ymin=0,ymax=8,ystep=2,ablesen]
\funktion{2*1.35^\x}{B}
\end{ksys}\ksysabstand\begin{ksys}[xmin=0,xmax=5,ymin=0,ymax=8,ystep=2,ablesen]
\gerade{0.8}{3}{C}
\end{ksys}

\begin{teile}
\teil bei Graph A
\teil bei Graph B
\teil bei Graph C
\teil Welcher Graph beginnt im Ursprung? \quad \kreuz{A}\kreuz{B}\kreuz{C}
\end{teile}
\end{aufgabe}

\verfahren{Wachstumsart an Tabelle und Text erkennen}

\begin{aufgabe}{Ich kann prüfen, ob in einer Tabelle immer derselbe Betrag dazukommt.}
Bilde die Differenzen benachbarter Werte. Kreuze an, ob sie gleich sind.
\begin{teile}
\teil 3; \ 7; \ 11; \ 15 \quad Differenzen: \leerfeld \quad gleich? \janein
\teil 2; \ 4; \ 8; \ 16 \quad Differenzen: \leerfeld \quad gleich? \janein
\teil 10; \ 25; \ 40; \ 55 \quad Differenzen: \leerfeld \quad gleich? \janein
\teil 5; \ 10; \ 20; \ 40 \quad Differenzen: \leerfeld \quad gleich? \janein
\teil 1,2; \ 1,8; \ 2,4; \ 3,0 \quad Differenzen: \leerfeld \quad gleich? \janein
\teil 90; \ 72; \ 54; \ 36 \quad Differenzen: \leerfeld \quad gleich? \janein
\end{teile}
\end{aufgabe}

\begin{aufgabe}{Ich kann prüfen, ob in einer Tabelle immer mit derselben Zahl multipliziert wird.}
Bilde die Quotienten: jeden Wert geteilt durch den Wert davor. Kreuze an, ob sie gleich sind.
\begin{teile}
\teil 3; \ 6; \ 12; \ 24 \quad Quotienten: \leerfeld \quad gleich? \janein
\teil 5; \ 15; \ 45; \ 135 \quad Quotienten: \leerfeld \quad gleich? \janein
\teil 2; \ 4; \ 6; \ 8 \quad Quotienten: \leerfeld \quad gleich? \janein
\teil 40; \ 50; \ 62,5; \ 78,125 \quad Quotienten: \leerfeld \quad gleich? \janein
\teil 800; \ 400; \ 200; \ 100 \quad Quotienten: \leerfeld \quad gleich? \janein
\teil 100; \ 110; \ 120; \ 130 \quad Quotienten: \leerfeld \quad gleich? \janein
\end{teile}
\end{aufgabe}

\begin{aufgabe}{Ich kann an einer Tabelle entscheiden, ob ein Wert linear oder exponentiell wächst.}
Entscheide und begründe in einem Satz.
\begin{teile}
\teil Kontostand nach 0, 1, 2, 3 Jahren: 200\,€; \ 260\,€; \ 320\,€; \ 380\,€
\teil Zahl der Follower nach 0, 1, 2, 3 Monaten: 500; \ 550; \ 605; \ 665,5
\teil Fläche eines Pilzgeflechts nach 0, 1, 2, 3 Wochen: 64\,cm²; \ 96\,cm²; \ 144\,cm²; \ 216\,cm²
\end{teile}
\end{aufgabe}

\begin{aufgabe}{Ich kann aus einem Text entscheiden, ob ein Wert linear oder exponentiell zu- oder abnimmt.}
Entscheide und begründe in einem Satz.
\begin{teile}
\teil Ein junger Baum wird jedes Jahr 40\,cm höher.
\teil Die Einwohnerzahl einer Stadt wächst jedes Jahr um 1,5\,\%.
\teil Ein Handy verliert jedes Jahr 25\,\% seines Werts.
\teil Aus einem Tank laufen jede Minute 12 Liter ab.
\teil Ein Fahrradverleih hatte diese Umsätze:

\sachtabelle{lcccc}{Jahr & 2021 & 2022 & 2023 & 2024}{Umsatz in € & 240\,000 & 252\,000 & 264\,600 & 277\,830}

Die Inhaberin sagt: „Der Umsatz wächst linear, denn er steigt jedes Jahr um 5\,\%.“ Entscheide, ob sie recht hat, und begründe. (P10 2018 OS)
\end{teile}
\end{aufgabe}

\verfahren{Graph zu einem Wachstum auswählen}

\begin{aufgabe}{Ich erkenne, ob ein Graph eine Gerade oder eine gekrümmte Kurve ist.}
Kreuze an.

\begin{ksys}[xmin=0,xmax=5,ymin=0,ymax=8,ystep=2,ablesen]
\gerade{1}{1}{A}
\end{ksys}\ksysabstand\begin{ksys}[xmin=0,xmax=5,ymin=0,ymax=8,ystep=2,ablesen]
\funktion{1.5^\x}{B}
\end{ksys}

\begin{ksys}[xmin=0,xmax=5,ymin=0,ymax=8,ystep=2,ablesen]
\funktion{0.5*2^\x}{C}
\end{ksys}\ksysabstand\begin{ksys}[xmin=0,xmax=5,ymin=0,ymax=8,ystep=2,ablesen]
\gerade{0.6}{2}{D}
\end{ksys}

\begin{teilezwei}
\tz A: \kreuz{Gerade}\kreuz{Kurve} & \tz B: \kreuz{Gerade}\kreuz{Kurve} \\
\tz C: \kreuz{Gerade}\kreuz{Kurve} & \tz D: \kreuz{Gerade}\kreuz{Kurve} \\
\end{teilezwei}
\end{aufgabe}

\begin{aufgabe}{Ich kann den passenden Graphen zu einem Wachstum auswählen.}
Auf der senkrechten Achse steht der Bestand in Tausend, auf der waagerechten die Zeit in Jahren. Kreuze an.

\begin{ksys}[xmin=0,xmax=5,ymin=0,ymax=8,ystep=2,ablesen]
\gerade{1.2}{0}{A}
\end{ksys}\ksysabstand\begin{ksys}[xmin=0,xmax=5,ymin=0,ymax=8,ystep=2,ablesen]
\funktion{2*1.4^\x}{B}
\end{ksys}\ksysabstand\begin{ksys}[xmin=0,xmax=5,ymin=0,ymax=8,ystep=2,ablesen]
\gerade{1.2}{2}{C}
\end{ksys}

\begin{teile}
\teil Ein Bestand von 2000 Tieren wächst jedes Jahr um 40\,\%. Welcher Graph kommt nicht in Frage, weil er nicht beim Anfangswert beginnt? \quad \kreuz{A}\kreuz{B}\kreuz{C}
\teil Welcher Graph passt zu diesem Bestand? \quad \kreuz{A}\kreuz{B}\kreuz{C}
\end{teile}

\begin{ksys}[xmin=0,xmax=5,ymin=0,ymax=8,ystep=2,ablesen]
\gerade{-1.4}{8}{D}
\end{ksys}\ksysabstand\begin{ksys}[xmin=0,xmax=5,ymin=0,ymax=8,ystep=2,ablesen]
\funktion{8*0.6^\x}{}
\node[fill=white,inner sep=1pt] at (1.7,5.4) {$E$};
\end{ksys}\ksysabstand\begin{ksys}[xmin=0,xmax=5,ymin=0,ymax=8,ystep=2,ablesen]
\funktion{1.5^\x}{F}
\end{ksys}

\begin{teile}
\teil Ein Bestand von 8000 nimmt jedes Jahr um 40\,\% ab. \quad \kreuz{D}\kreuz{E}\kreuz{F}
\teil Ein Bestand von 8000 nimmt jedes Jahr um 1400 ab. \quad \kreuz{D}\kreuz{E}\kreuz{F}
\end{teile}
\end{aufgabe}

\begin{aufgabe}{Ich kann begründen, warum ein Graph nicht zu einem Wachstum passt.}
Die Zahl der Nutzer einer App beginnt bei 3000 und wächst jeden Monat um 25\,\%. Senkrechte Achse: Nutzer in Tausend, waagerechte Achse: Zeit in Monaten. Begründe in einem Satz.

\begin{ksys}[xmin=0,xmax=5,ymin=0,ymax=8,ystep=2,ablesen]
\gerade{0.9}{3}{G}
\end{ksys}\ksysabstand\begin{ksys}[xmin=0,xmax=5,ymin=0,ymax=8,ystep=2,ablesen]
\funktion{0.35*\x^2}{H}
\end{ksys}\ksysabstand\begin{ksys}[xmin=0,xmax=5,ymin=0,ymax=8,ystep=2,ablesen]
\funktion{3*1.25^\x}{I}
\end{ksys}

\begin{teile}
\teil Warum passt der Graph G nicht?
\teil Warum passt der Graph H nicht?
\end{teile}

In einem See leben 6000 Fische. Wegen einer Krankheit nimmt ihre Zahl jedes Jahr um 30\,\% ab. Senkrechte Achse: Fische in Tausend, waagerechte Achse: Zeit in Jahren.

\begin{ksys}[xmin=0,xmax=5,ymin=0,ymax=8,ystep=2,ablesen]
\funktion{6*0.7^\x}{}
\node[fill=white,inner sep=1pt] at (1.7,5) {$J$};
\end{ksys}\ksysabstand\begin{ksys}[xmin=0,xmax=5,ymin=0,ymax=8,ystep=2,ablesen]
\gerade{-1.2}{6}{K}
\end{ksys}\ksysabstand\begin{ksys}[xmin=0,xmax=5,ymin=0,ymax=8,ystep=2,ablesen]
\funktion{8*0.7^\x}{L}
\end{ksys}

\begin{teile}
\teil Gib an, welcher Graph passt, und begründe für jeden der beiden anderen Graphen, warum er nicht passt. (P10 2026 FOR)
\end{teile}
\end{aufgabe}

\verfahren{Wertepaare als Punkte darstellen}

\begin{aufgabe}{Ich kann Wertepaare eines Wachstums als Punkte eintragen.}
Eine Zellkultur wird beobachtet ($t$: Zeit in Stunden, $N$: Zahl der Zellen). Trage die Punkte aus der Tabelle ein.

\wertetabelle[100,200,400]{t}{N}{0,1,2}

\begin{teile}
\teil Punkt für $t = 0$
\teil Punkt für $t = 1$
\teil Punkt für $t = 2$
\end{teile}

\begin{ksys}[xmin=0,xmax=4,ymin=0,ymax=500,ystep=100,xlabel=$t$ in h,ylabel=$N$]
\end{ksys}

\end{aufgabe}

\begin{aufgabe}{Ich kann Wertepaare eines Wachstums als Punkte eintragen – weiter.}
Eine zweite Kultur ($t$ in Stunden, $N$ Zahl der Zellen). Die ersten beiden Punkte sind schon eingetragen.

\wertetabelle[100,150,225,{337{,}5}]{t}{N}{0,1,2,3}

\begin{teile}
\teil Trage den Punkt für $t = 2$ genau ein.
\teil Trage den Punkt für $t = 3$ genau ein.
\teil Verbinde die vier Punkte zu einer Kurve.
\end{teile}

\begin{ksys}[xmin=0,xmax=4,ymin=0,ymax=400,ystep=50,xlabel=$t$ in h,ylabel=$N$]
\punkt{0}{100}{}
\punkt{1}{150}{}
\end{ksys}

\end{aufgabe}

\begin{aufgabe}{Ich kann für Wertepaare eine passende Achseneinteilung wählen.}
Ein E-Bike verliert jedes Jahr 20\,\% seines Werts.

\sachtabelle{lccccc}{Zeit in Jahren & 0 & 1 & 2 & 3 & 4}{Wert in € & 2000 & 1600 & 1280 & 1024 & 819,20}

\begin{teile}
\teil Teile beide Achsen passend ein und beschrifte sie.
\teil Trage die fünf Wertepaare als Punkte ein.
\teil Verbinde die Punkte zu einer Kurve. (P10 2017 OS)
\end{teile}

\leeresgitter{10}{10}

\end{aufgabe}

\begin{aufgabe}{Ich finde den Fehler bei der Entscheidung zwischen linear und exponentiell.}
Finde den Fehler, benenne ihn und korrigiere die Aussage.
\begin{teile}
\teil Ein Verein hat nacheinander 300, 330, 363 und 399,3 Mitglieder. Mia sagt: „Das Wachstum ist linear, denn jedes Jahr kommen 10\,\% dazu.“
\teil Ein Guthaben beginnt bei 500\,€ und wächst jedes Jahr um 2\,\%. Ben zeichnet dazu eine Gerade durch den Ursprung, „weil jedes Jahr gleich viel Prozent dazukommen“.
\end{teile}
\end{aufgabe}

\begin{aufgabe}{Ich kann begründen, warum bei gleichem Prozentsatz der Zuwachs immer größer wird.}
Begründe.
\begin{teile}
\teil Ein Guthaben von 1000\,€ wächst jedes Jahr um 10\,\%. Zeige an den ersten beiden Jahren, dass der Zuwachs in Euro größer wird.
\teil Warum heißt „jedes Jahr gleich viel Prozent“ nicht „jedes Jahr gleich viel Euro“?
\end{teile}
\end{aufgabe}

\begin{aufgabe}{Ich kann Grenzen eines Wachstumsmodells nennen.}
Begründe.
\begin{teile}
\teil Eine Sonnenblume ist 25\,cm hoch und wächst in den ersten Wochen jede Woche um 40\,\%. Kann sie ein ganzes Jahr so weiterwachsen?
\teil In einer Schule mit 800 Schülern verdoppelt sich jeden Tag die Zahl derer, die ein Gerücht kennen. Warum kann das nicht lange so weitergehen?
\end{teile}
\end{aufgabe}

\begin{aufgabe}{Ich kann zwei Sparangebote mit linearem und exponentiellem Wachstum vergleichen.}
Konto A: 1000\,€, jedes Jahr kommen 60\,€ dazu. \quad Konto B: 1000\,€, jedes Jahr kommen 5\,\% Zinsen dazu, die mitverzinst werden.
\begin{teile}
\teil Berechne beide Kontostände nach 1, 2 und 3 Jahren.
\teil Rechne weiter: Nach welchem Jahr hat Konto B zum ersten Mal mehr Geld als Konto A?
\teil Welches Konto ist besser, wenn du 5 Jahre sparst, welches bei 10 Jahren? Begründe mit den Kontoständen.
\end{teile}
\end{aufgabe}

\clearpage
% Einheit 2 – Nr. 30–39
\setcounter{aufgabe}{29}
\einheitenkopf[e2][2 Wachstumsfaktor und Tabelle]{Einheit 2 von 5 · Wachstumsfaktor und Wachstumstabelle}
\zweigzeile{Hier lernst du, aus einem Prozentsatz den Wachstumsfaktor zu bilden und Wachstumstabellen zu ergänzen · neu in diesem Jahr · P10 oft · baut auf: Differenz oder Quotient (Einheit 1), Erhöhen um p\,\%, Quotient zweier Werte}

\begin{aufgabe}{Ich erkenne am Faktor, ob ein Wert zunimmt oder abnimmt.}
Kreuze an. Rechne nichts.
\begin{teile}
\teil Der Wert wird mit 1,4 multipliziert. \quad \kreuz{nimmt zu}\kreuz{nimmt ab}
\teil Der Wert wird mit 0,7 multipliziert. \quad \kreuz{nimmt zu}\kreuz{nimmt ab}
\teil Der Wert wird mit 1,05 multipliziert. \quad \kreuz{nimmt zu}\kreuz{nimmt ab}
\teil Der Wert wird mit 0,95 multipliziert. \quad \kreuz{nimmt zu}\kreuz{nimmt ab}
\teil Der Wert wird mit 1,002 multipliziert. \quad \kreuz{nimmt zu}\kreuz{nimmt ab}
\end{teile}
\end{aufgabe}

\verfahren{Wachstumsfaktor bestimmen}

\begin{aufgabe}{Ich kann aus einem Prozentsatz den Wachstumsfaktor bilden.}
Gib den Wachstumsfaktor an.
\begin{teilezwei}
\tz Zunahme um 10\,\% \quad \feld{q} & \tz Zunahme um 20\,\% \quad \feld{q} \\
\tz Zunahme um 50\,\% \quad \feld{q} & \tz Zunahme um 3\,\% \quad \feld{q} \\
\tz Zunahme um 2,5\,\% \quad \feld{q} & \tz Abnahme um 10\,\% \quad \feld{q} \\
\tz Abnahme um 25\,\% \quad \feld{q} & \tz Abnahme um 4\,\% \quad \feld{q} \\
\tz Abnahme um 0,5\,\% \quad \feld{q} & \tz Zunahme um 100\,\% \quad \feld{q} \\
\end{teilezwei}
\end{aufgabe}

\begin{aufgabe}{Ich kann aus einem Wachstumsfaktor die prozentuale Änderung ablesen.}
Gib an, um wie viel Prozent der Wert zu- oder abnimmt. Schreibe + für Zunahme und $-$ für Abnahme.
\begin{teilezwei}
\tz $q = 1{,}04$ \quad \leerfeld[\%] & \tz $q = 1{,}25$ \quad \leerfeld[\%] \\
\tz $q = 1{,}08$ \quad \leerfeld[\%] & \tz $q = 1{,}6$ \quad \leerfeld[\%] \\
\tz $q = 0{,}9$ \quad \leerfeld[\%] & \tz $q = 0{,}85$ \quad \leerfeld[\%] \\
\tz $q = 1{,}015$ \quad \leerfeld[\%] & \tz $q = 0{,}995$ \quad \leerfeld[\%] \\
\tz $q = 3$ \quad \leerfeld[\%] & \\
\end{teilezwei}
\end{aufgabe}

\begin{aufgabe}{Ich kann den Wachstumsfaktor als Quotient aus einer Tabelle bestimmen.}
Berechne die Quotienten benachbarter Werte und gib den Wachstumsfaktor an.
\begin{teile}
\teil 20; \ 40; \ 80 \quad \feld{q}
\teil 7; \ 21; \ 63 \quad \feld{q}
\teil 300; \ 330; \ 363 \quad \feld{q}
\teil 500; \ 400; \ 320 \quad \feld{q}
\teil 1200; \ 1260; \ 1323; \ 1389,15 \quad Sind alle Quotienten gleich? Antworte in einem Satz und gib $q$ an.
\teil Ein Bestand ist im Jahr 0 genau 400 und im Jahr 2 genau 484. Er wächst jedes Jahr um denselben Prozentsatz. Bestimme den Wachstumsfaktor und den Prozentsatz.
\teil Die Zahl der Bakterien in einer Probe wurde stündlich gezählt:

\sachtabelle{lcccc}{Zeit in h & 0 & 1 & 2 & 3}{Anzahl & 128 & 160 & 200 & 250}

Weise nach, dass die Anzahl jede Stunde mit demselben Faktor wächst, und gib an, um wie viel Prozent sie stündlich zunimmt. (P10 2025 OS)
\end{teile}
\end{aufgabe}

\verfahren{Wachstumstabelle fortschreiben}

\begin{aufgabe}{Ich kann eine Wachstumstabelle mit dem Faktor fortschreiben.}
Ergänze die Tabelle.
\begin{teile}
\teil Wachstumsfaktor 2 \\ \wertetabelle[5,,,]{t}{N}{0,1,2,3}
\teil Wachstumsfaktor 3 \\ \wertetabelle[2,,,]{t}{N}{0,1,2,3}
\teil Wachstumsfaktor 1,1 \\ \wertetabelle[100,,,]{t}{N}{0,1,2,3}
\teil Wachstumsfaktor 1,5 \\ \wertetabelle[40,,,]{t}{N}{0,1,2,3}
\end{teile}
\end{aufgabe}

\begin{aufgabe}{Ich kann eine Wachstumstabelle mit dem Faktor fortschreiben – weiter.}
Ergänze die Tabelle.
\begin{teile}
\teil Ein Bestand von 250 Fischen wächst jedes Jahr um 20\,\%. Trage den Anfangswert und die Werte der nächsten beiden Jahre ein. \\ \wertetabelle[,,]{t}{N}{0,1,2}
\teil Wachstumsfaktor 1,2 \\ \wertetabelle[150,,216,]{t}{N}{0,1,2,3}
\teil Ein Guthaben von 800\,€ wächst jedes Jahr um 5\,\%. Berechne den Wert nach sechs Jahren in einem Schritt. \\ \wertetabelle[800,]{t}{K}{0,6}
\end{teile}
\end{aufgabe}

\begin{aufgabe}{Ich kann eine Wachstumstabelle mit dem Faktor fortschreiben – weiter.}
Ergänze die Tabelle.
\begin{teile}
\teil Wachstumsfaktor 1,1. Berechne den Wert davor. \\ \wertetabelle[,440,484]{t}{N}{0,1,2}
\teil Wachstumsfaktor 1,5. Bestimme die fehlende Zeitangabe und prüfe sie mit dem Faktor. \\ \wertetabelle[128,192,288,972]{t}{N}{0,1,2,{\;}}
\teil Wachstumsfaktor 0,75 \\ \wertetabelle[2400,,,]{t}{N}{0,1,2,3}
\teil Ein Medikament wird im Körper abgebaut, jede Stunde um 15\,\%. Ergänze den fehlenden Wert und die fehlende Zeitangabe. (P10 2016 OS) \\ \wertetabelle[40,34,,{20{,}88}]{t \text{ in h}}{m \text{ in mg}}{0,1,2,{\;}}
\end{teile}
\end{aufgabe}

\begin{aufgabe}{Ich finde den Fehler in einer Wachstumstabelle.}
Finde den Fehler, benenne ihn und korrigiere die Tabelle.
\begin{teile}
\teil Ein Bestand von 500 wächst jährlich um 4\,\%. Lena schreibt: \\ \wertetabelle[500,520,540,560]{t}{N}{0,1,2,3}
\teil Ein Bestand von 200 wächst jährlich um 5\,\%. Kai rechnet mit dem Faktor 0,05: \\ \wertetabelle[200,10,{0{,}5}]{t}{N}{0,1,2}
\teil Ein Guthaben von 800\,€ wächst jährlich um 3\,\%. Jonas schreibt: \\ \wertetabelle[824,{848{,}72}]{t}{K}{0,1}
\end{teile}
\end{aufgabe}

\begin{aufgabe}{Ich kann begründen, woran man den Wachstumsfaktor erkennt.}
Begründe.
\begin{teile}
\teil Warum gehört zu „plus 8\,\%“ der Faktor 1,08 und nicht 0,08?
\teil Ein Bestand hat die Werte 200; \ 220; \ 242; \ 266,2. Tim sagt: „Die Differenzen 20; 22; 24,2 werden größer, also gibt es keinen festen Faktor.“ Hat er recht?
\end{teile}
\end{aufgabe}

\begin{aufgabe}{Ich kann mit einer Wachstumstabelle eine Frage im Sachzusammenhang beantworten.}
Seerosen bedecken 12\,m² eines 50\,m² großen Teichs. Die bedeckte Fläche wächst jede Woche um 25\,\%.
\begin{teile}
\teil Lege eine Tabelle für die Wochen 0 bis 6 an.
\teil In welcher Woche ist zum ersten Mal mehr als die Hälfte des Teichs bedeckt?
\teil Der Besitzer will Seerosen entfernen, bevor 40\,m² bedeckt sind. Bis zu welcher Woche muss er es spätestens tun? Begründe mit der Tabelle.
\end{teile}
\end{aufgabe}

\clearpage
% Einheit 3 – Nr. 40–49
\setcounter{aufgabe}{39}
\einheitenkopf[e3][3 Exponentialfunktion]{Einheit 3 von 5 · Exponentialfunktion aufstellen und auswerten}
\zweigzeile{Hier lernst du, eine Wachstumsgleichung $N(t) = N_0 \cdot q^t$ aufzustellen, zu deuten und damit Werte und Schwellen zu berechnen · neu in diesem Jahr · P10 · baut auf: Wachstumsfaktor und Tabelle (Einheit 2), Potenz mit dem Taschenrechner}

\begin{aufgabe}{Ich erkenne, welches Jahr der Schritt null ist.}
Unterstreiche das Startjahr und kreise das gesuchte Jahr ein. Rechne nichts.
\begin{teile}
\teil 2020 hatte ein Verein 150 Mitglieder. Die Zahl wächst jährlich um 5\,\%. Wie viele Mitglieder hat er 2025?
\teil Im Jahr 2018 lebten 900 Störche in einer Region, ihre Zahl nimmt jährlich um 2\,\% ab. Wie viele sind es im Jahr 2030?
\teil Wie groß ist im Jahr 2031 ein Guthaben, das 2024 bei 1500\,€ lag und seitdem jährlich mit 3\,\% verzinst wird?
\teil Ein Auto kostete 2022 neu 30\,000\,€. Welchen Wert hat es 2027, wenn es jährlich 15\,\% seines Werts verliert?
\end{teile}
\end{aufgabe}

\verfahren{Wachstumsgleichung aufstellen und deuten}

\begin{aufgabe}{Ich kann eine Wachstumsgleichung aufstellen.}
Gib den Anfangswert und den Wachstumsfaktor an und stelle die Gleichung auf.
\begin{teile}
\teil In einer Probe sind 100 Bakterien, ihre Zahl verdoppelt sich jede Stunde. \\ \feld{N_0} \quad \feld{q} \quad $N(t) = $ \leerfeld
\teil Auf einem Teich schwimmen 50 Wasserlinsen, ihre Zahl verdreifacht sich jede Woche. \\ \feld{N_0} \quad \feld{q} \quad $N(t) = $ \leerfeld
\teil Ein Guthaben von 400\,€ wächst jährlich um 10\,\%. \\ \feld{N_0} \quad \feld{q} \quad $N(t) = $ \leerfeld
\teil Eine Stadt mit 2500 Einwohnern wächst jährlich um 2\,\%. \\ \feld{N_0} \quad \feld{q} \quad $N(t) = $ \leerfeld
\teil Im Blut sind 90\,mg eines Wirkstoffs, stündlich werden 20\,\% abgebaut. \\ \feld{N_0} \quad \feld{q} \quad $N(t) = $ \leerfeld
\teil Ein Gletscher bedeckt 12\,km², die Fläche nimmt jährlich um 1,5\,\% ab. \\ \feld{N_0} \quad \feld{q} \quad $N(t) = $ \leerfeld
\end{teile}
\end{aufgabe}

\begin{aufgabe}{Ich kann die passende Gleichung unter mehreren auswählen.}
Kreuze die passende Gleichung an. Dabei ist $x$ die Zeit in Jahren und $y$ der Bestand.
\begin{teile}
\teil Ein Bestand von 50 Tieren verdoppelt sich jedes Jahr. \\ \kreuz{$y = 50 \cdot 2^x$}\quad\kreuz{$y = 2 \cdot 50^x$}\quad\kreuz{$y = 50 \cdot x^2$}\quad\kreuz{$y = 50 + 2x$}
\teil Ein Bestand von 600 Tieren nimmt jährlich um 5\,\% ab. \\ \kreuz{$y = 600 \cdot 1{,}05^x$}\quad\kreuz{$y = 600 \cdot 0{,}95^x$}\quad\kreuz{$y = 600 \cdot 0{,}5^x$}\quad\kreuz{$y = 0{,}95 \cdot 600^x$}
\teil Ein Kapital von 3000\,€ wächst jährlich um 1,2\,\%. (P10 2017 OS) \\ \kreuz{$y = 3000 \cdot 1{,}2^x$}\quad\kreuz{$y = 3000 \cdot 1{,}012^x$}\quad\kreuz{$y = 3000 \cdot 0{,}012^x$}\quad\kreuz{$y = 3000 + 1{,}012x$}
\end{teile}
\end{aufgabe}

\begin{aufgabe}{Ich kann die Teile einer Wachstumsgleichung im Sachzusammenhang deuten.}
Gib an, was die Zahlen und die Variable bedeuten.
\begin{teile}
\teil $N(t) = 300 \cdot 2^t$ beschreibt die Zahl der Hefezellen in einer Probe, $t$ in Stunden.
\teil $W(t) = 18\,000 \cdot 0{,}85^t$ beschreibt den Wert eines Autos in Euro, $t$ in Jahren.
\teil Erkläre, warum der Faktor 1,04 für eine Zunahme um 4\,\% steht.
\teil $f(x) = 50 \cdot 1{,}03^x$ gibt die Masse eines Fohlens in kg an, $x$ in Wochen nach der Geburt. Gib die Bedeutung von 50, von 1,03 und von $x$ an. (P10 2019 OS)
\end{teile}
\end{aufgabe}

\verfahren{Werte und Schwellenwerte berechnen}

\begin{aufgabe}{Ich kann mit einer Wachstumsgleichung den Wert nach einigen Schritten berechnen.}
Berechne den Wert.
\begin{teilezwei}
\tz $N(t) = 5 \cdot 2^t$; \ $N(3)$ \leerfeld & \tz $N(t) = 10 \cdot 3^t$; \ $N(2)$ \leerfeld \\
\tz $N(t) = 100 \cdot 1{,}1^t$; \ $N(2)$ \leerfeld & \tz $N(t) = 1000 \cdot 0{,}5^t$; \ $N(3)$ \leerfeld \\
\end{teilezwei}
\anweisung{Rechne mit dem Taschenrechner und runde erst am Ende sinnvoll.}
\begin{teile}
\teil $N(t) = 250 \cdot 1{,}08^t$; \ $N(10)$ \leerfeld
\teil Eine Stadt hatte im Jahr 2021 genau 4000 Einwohner, die Zahl wächst jährlich um 1,5\,\%. Wie viele Einwohner hat sie im Jahr 2029? \leerfeld
\teil Ein Auto kostet neu 24\,000\,€ und verliert jährlich 12\,\% seines Werts. Wie viel ist es nach fünf Jahren noch wert? \leerfeld[€]
\end{teile}
\end{aufgabe}

\begin{aufgabe}{Ich kann eine Behauptung über ein Wachstum mit einer Rechnung prüfen.}
Prüfe mit einer Rechnung und entscheide.
\begin{teile}
\teil $N(t) = 300 \cdot 1{,}5^t$. Behauptung: Nach vier Schritten ist der Bestand größer als 1500.
\teil Eine Gemeinde hat 5000 Einwohner und wächst jährlich um 2\,\%. Studie A sagt: „In zehn Jahren sind es 20\,\% mehr, also 6000.“ Studie B rechnet mit dem Wachstumsfaktor. Welche Studie sagt mehr voraus, und um wie viele Einwohner?
\teil Algen bedecken 30\,m² eines Sees, die Fläche wächst täglich um 25\,\%. Ein Zeitungsartikel behauptet: „Nach drei Wochen bedecken die Algen mehr als 3000\,m².“ Stelle eine Gleichung auf und prüfe die Behauptung. (P10 2025 OS)
\end{teile}
\end{aufgabe}

\begin{aufgabe}{Ich kann berechnen, wann ein Wert eine Schwelle überschreitet.}
Bestimme, nach wie vielen Schritten die Schwelle zum ersten Mal überschritten ist, und gib den Wert dann an.
\begin{teile}
\teil $N(t) = 10 \cdot 2^t$, Schwelle 100 \quad \feld{t} \feld{N}
\teil $N(t) = 3 \cdot 3^t$, Schwelle 200 \quad \feld{t} \feld{N}
\teil $N(t) = 500 \cdot 1{,}25^t$, Schwelle 1000 \quad \feld{t} \feld{N}
\end{teile}
\anweisung{Bestimme, nach wie vielen Schritten die Schwelle zum ersten Mal unterschritten ist.}
\begin{teile}
\teil $N(t) = 800 \cdot 0{,}7^t$, Schwelle 100 \quad \feld{t} \feld{N}
\end{teile}
\anweisung{Beantworte die Frage.}
\begin{teile}
\teil Ein Start-up hatte im Jahr 2022 einen Umsatz von 420\,000\,€ und steigert ihn jährlich um 15\,\%. In welchem Jahr übersteigt der Umsatz zum ersten Mal eine Million Euro? Gib den Umsatz in diesem Jahr an. (P10 2018 OS)
\end{teile}
\end{aufgabe}

\clearpage
\begin{aufgabe}{Ich finde den Fehler in einer Rechnung mit der Wachstumsgleichung.}
Finde den Fehler, benenne ihn und korrigiere die Rechnung.
\begin{teile}
\teil Ein Bestand von 600 Tieren wächst seit 2020 jährlich um 3\,\%. Jonas berechnet den Bestand im Jahr 2026:
\rechnung{N(t) &= 600 \cdot 1{,}03^t \\ N(2026) &= 600 \cdot 1{,}03^{2026}}
\teil Ein Guthaben von 1200\,€ wächst jährlich um 4\,\%. Sara berechnet das Guthaben nach fünf Jahren:
\rechnung{K &= 1200\,€ \cdot 1{,}04 \cdot 5 \\ K &= 6240\,€}
\teil Ein Bestand von 90 Tieren wächst jährlich um 4\,\%. Tim stellt die Gleichung auf:
\rechnung{N(t) &= 90 \cdot 4^t}
\end{teile}
\end{aufgabe}

\begin{aufgabe}{Ich kann begründen, warum in der Wachstumsgleichung eine Potenz steht.}
Begründe.
\begin{teile}
\teil Warum rechnet man bei einem Wachstum um 10\,\% in drei Jahren $500 \cdot 1{,}1^3$ und nicht $500 \cdot 1{,}1 \cdot 3$?
\teil Warum ist „dreimal nacheinander um 10\,\% erhöht“ mehr als „einmal um 30\,\% erhöht“?
\end{teile}
\end{aufgabe}

\begin{aufgabe}{Ich kann mit Wachstumsgleichungen zwei Entwicklungen vergleichen.}
Stadt A hat 40\,000 Einwohner und wächst jährlich um 1,5\,\%. Stadt B hat 46\,000 Einwohner und schrumpft jährlich um 1\,\%.
\begin{teile}
\teil Stelle für beide Städte eine Gleichung auf.
\teil Berechne beide Einwohnerzahlen nach fünf Jahren.
\teil Nach wie vielen Jahren hat Stadt A zum ersten Mal mehr Einwohner als Stadt B?
\end{teile}
\end{aufgabe}

\clearpage
% Einheit 4 – Nr. 50–57
\setcounter{aufgabe}{49}
\einheitenkopf[e4][4 Verdopplungs- und Halbwertszeit]{Einheit 4 von 5 · Verdopplungs- und Halbwertszeit}
\zweigzeile{Hier lernst du, die Zeit zu bestimmen, nach der sich ein Wert verdoppelt oder halbiert hat · neu in diesem Jahr · P10 · baut auf: Wachstumstabelle (Einheit 2), Wachstumsgleichung (Einheit 3), Werte am Graphen ablesen}

\begin{aufgabe}{Ich erkenne, ob sich die Zeit oder der Bestand verdoppeln oder halbieren soll.}
Kreuze an, was sich verdoppeln oder halbieren soll. Rechne nichts.
\begin{teile}
\teil Nach wie vielen Jahren hat sich das Guthaben verdoppelt? \quad \kreuz{die Zeit}\kreuz{der Bestand}
\teil Wie lange dauert es, bis nur noch die Hälfte des Wirkstoffs im Blut ist? \quad \kreuz{die Zeit}\kreuz{der Bestand}
\teil Wie viele Tiere sind es, wenn man doppelt so lange wartet? \quad \kreuz{die Zeit}\kreuz{der Bestand}
\teil Nach welcher Zeit sind aus 300 Bakterien 600 geworden? \quad \kreuz{die Zeit}\kreuz{der Bestand}
\teil Wie viel ist nach der halben Beobachtungszeit noch da? \quad \kreuz{die Zeit}\kreuz{der Bestand}
\end{teile}
\end{aufgabe}

\begin{aufgabe}{Ich kann den Zielwert für die Verdopplung oder Halbierung angeben.}
Gib den Wert an, der erreicht sein muss.
\begin{teile}
\teil Verdopplung, Start bei 40 \quad \leerfeld
\teil Verdopplung, Start bei 250 \quad \leerfeld
\teil Halbierung, Start bei 90 \quad \leerfeld
\teil Halbierung, Start bei 3,6\,g \quad \leerfeld[g]
\teil Ein Guthaben startet bei 500\,€ und hat nach fünf Jahren 700\,€. Welchen Betrag muss es für die Verdopplung erreichen? \leerfeld[€]
\end{teile}
\end{aufgabe}

\begin{aufgabe}{Ich kann die Verdopplungs- oder Halbwertszeit aus einer Tabelle ablesen.}
Lies die Verdopplungszeit ab.
\begin{teile}
\teil \wertetabelle[5,10,20,40]{t \text{ in h}}{N}{0,1,2,3} \quad \feld[h]{T}
\end{teile}
\anweisung{Der Zielwert liegt zwischen zwei Tabellenwerten. Gib die nächstliegende Zeit an.}
\begin{teile}
\teil Verdopplungszeit: \\ \wertetabelle[16,24,36,54,81]{t \text{ in Jahren}}{N}{0,1,2,3,4} \quad \feld[Jahre]{T}
\teil Halbwertszeit: \\ \wertetabelle[80,56,{39{,}2},{27{,}44}]{t \text{ in h}}{m \text{ in mg}}{0,1,2,3} \quad \feld[h]{T}
\teil Ein radioaktives Präparat zerfällt. Gib seine Halbwertszeit an. (P10 2016 OS) \\ \wertetabelle[12,{9{,}6},{7{,}68},{6{,}14},{4{,}92}]{t \text{ in Tagen}}{m \text{ in mg}}{0,2,4,6,8} \quad \feld[Tage]{T}
\end{teile}
\end{aufgabe}

\begin{aufgabe}{Ich kann die Verdopplungs- oder Halbwertszeit am Graphen ablesen.}
Links: Wirkstoff im Blut ($t$ in Stunden, $m$ in mg). Rechts: Zahl der Kaninchen auf einer Insel ($t$ in Monaten).

\begin{ksys}[xmin=0,xmax=8,ymin=0,ymax=400,ystep=50,ablesen,xlabel=$t$,ylabel=$m$]
\funktion{400*0.8^\x}{}
\end{ksys}\ksysabstand\begin{ksys}[xmin=0,xmax=10,ymin=0,ymax=180,ystep=20,ablesen,xlabel=$t$,ylabel=$N$]
\funktion{60*1.12^\x}{}
\end{ksys}

\begin{teile}
\teil Links: Lies ab, nach welcher Zeit nur noch die Hälfte des Wirkstoffs im Blut ist. \feld[h]{T}
\teil Beschreibe in zwei Sätzen, wie du bei a) vorgegangen bist.
\teil Rechts: Ermittle mithilfe des Graphen die Verdopplungszeit der Kaninchenzahl und beschreibe dein Vorgehen. (P10 2020 OS)
\end{teile}
\end{aufgabe}

\begin{aufgabe}{Ich kann die Verdopplungs- oder Halbwertszeit durch Probieren mit dem Faktor bestimmen.}
Bestimme durch Probieren, nach wie vielen ganzen Schritten sich der Wert ungefähr verdoppelt oder halbiert hat.
\begin{teile}
\teil Faktor 1,5, Verdopplung \quad \leerfeld
\teil Faktor 1,25, Verdopplung \quad \leerfeld
\teil Faktor 0,5, Halbierung \quad \leerfeld
\teil Faktor 0,8, Halbierung \quad \leerfeld
\teil Zunahme um 10\,\% je Schritt, Verdopplung \quad \leerfeld
\teil Abnahme um 10\,\% je Schritt, Halbierung \quad \leerfeld
\anweisung{Prüfe mit einer Rechnung.}
\teil Ein Wert wächst jährlich um 10\,\%. Berechne für die Startwerte 80 und 500 den Wert nach sieben Jahren. Was fällt dir auf?
\end{teile}
\anweisung{Löse mit dem Logarithmus (Taste log).}
\begin{teile}
\teil Bestimme die Verdopplungszeit bei 10\,\% Zunahme pro Jahr auf zwei Nachkommastellen genau. \feld[Jahre]{T}
\end{teile}
\end{aufgabe}

\begin{aufgabe}{Ich finde den Fehler bei der Verdopplungs- oder Halbwertszeit.}
Finde den Fehler, benenne ihn und korrigiere.
\begin{teile}
\teil Ein Bestand wächst. Paul sagt: „Die Verdopplungszeit ist 2 Jahre, denn 2 Jahre sind das Doppelte von 1 Jahr.“ \\ \wertetabelle[300,360,432,{518{,}4},{622{,}08}]{t \text{ in Jahren}}{N}{0,1,2,3,4}
\teil Lea soll am Graphen die Halbwertszeit eines Stoffes ablesen, der mit 8\,mg beginnt. Sie sucht die 4 auf der waagerechten Achse, geht senkrecht bis zum Graphen und liest an der senkrechten Achse den Wert ab.
\teil Ein Stoff zerfällt. Maja sagt: „Die Halbwertszeit ist 3 Stunden, weil 29,64 der letzte Wert über 25 ist.“ \\ \wertetabelle[50,42,{35{,}28},{29{,}64},{24{,}89}]{t \text{ in h}}{m \text{ in mg}}{0,1,2,3,4}
\end{teile}
\end{aufgabe}

\begin{aufgabe}{Ich kann begründen, wann es eine feste Verdopplungszeit gibt.}
Begründe.
\begin{teile}
\teil Ein Wert wächst jedes Jahr um 10\,\%. Warum dauert es von 100 bis 200 genauso lange wie von 1000 bis 2000?
\teil Ein Wert wächst linear: 100; \ 120; \ 140; \ 160; \ … \ Gibt es eine feste Verdopplungszeit?
\end{teile}
\end{aufgabe}

\begin{aufgabe}{Ich kann mit der Halbwertszeit eine Frage im Sachzusammenhang beantworten.}
Ein Schmerzmittel hat im Blut eine Halbwertszeit von 6 Stunden. Um 8 Uhr sind 240\,mg im Blut.
\begin{teile}
\teil Wie viel ist um 14 Uhr, um 20 Uhr und um 2 Uhr nachts noch im Blut?
\teil Eine neue Tablette darf erst genommen werden, wenn weniger als 70\,mg im Blut sind. Zu welcher der Uhrzeiten aus a) ist das zuerst der Fall?
\teil Eine Freundin sagt: „Nach 24 Stunden ist nichts mehr im Blut.“ Stimmt das? Begründe mit einer Rechnung.
\end{teile}
\end{aufgabe}

\clearpage
% Einheit 5 – Nr. 58–69
\setcounter{aufgabe}{57}
\einheitenkopf[e5][5 Potenzfunktionen]{Einheit 5 von 5 · Potenzfunktionen mit natürlichem Exponenten}
\zweigzeile{Hier lernst du Potenzfunktionen wie $y = x^3$ und $y = x^4$ kennen: Wertetabelle, Graph, Symmetrie und Funktionswerte · neu in diesem Jahr (am Gymnasium schon seit Klasse 9) · keine P10-Aufgabe · baut auf: Potenz mit dem Taschenrechner, Punkte eintragen, Normalparabel}

\verfahren{Wertetabelle und Graph}

\begin{aufgabe}{Ich kann eine Wertetabelle zu einer Potenzfunktion ausfüllen.}
Ergänze die Wertetabelle.
\begin{teile}
\teil $y = x^3$ \\ \wertetabelle{x}{y}{0,1,2,3}
\teil $y = x^4$ \\ \wertetabelle{x}{y}{0,1,2,3}
\teil $y = x^3$ \\ \wertetabelle{x}{y}{4,5,10}
\teil $y = x^4$ \\ \wertetabelle{x}{y}{4,5,10}
\teil $y = x^3$ \\ \wertetabelle{x}{y}{-3,-2,-1}
\teil $y = x^4$ \\ \wertetabelle{x}{y}{-3,-2,-1}
\end{teile}
\end{aufgabe}

\begin{aufgabe}{Ich kann eine Wertetabelle zu einer Potenzfunktion ausfüllen – weiter.}
Ergänze die Wertetabelle.
\begin{teile}
\teil $y = 0{,}5x^3$ \\ \wertetabelle{x}{y}{-2,-1,0,1,2}
\teil $y = -x^4$ \\ \wertetabelle{x}{y}{-2,-1,0,1,2}
\teil $y = x^3$ \\ \wertetabelle{x}{y}{{-1{,}5},{-0{,}5},{0{,}5},{1{,}5}}
\end{teile}
\end{aufgabe}

\begin{aufgabe}{Ich kann den Graphen einer Potenzfunktion zeichnen.}
Zeichne die Graphen mithilfe deiner Wertetabellen aus Nr. 58 und 59 in das Koordinatensystem und beschrifte sie.
\begin{teile}
\teil $y = x^3$
\teil $y = x^4$
\teil $y = 0{,}5x^3$
\end{teile}

\begin{ksys}[xmin=-2,xmax=2,ymin=-8,ymax=8,xstep=0.5,ystep=2]
\end{ksys}

\end{aufgabe}

\verfahren{Funktionswerte, Punktprobe und Umkehrung}

\begin{aufgabe}{Ich kann Funktionswerte einer Potenzfunktion berechnen.}
Es ist $f(x) = x^3$, $g(x) = x^4$ und $h(x) = 0{,}5x^3$. Berechne.
\begin{teilezwei}
\tz $f(2)$ \leerfeld & \tz $f(4)$ \leerfeld \\
\tz $g(3)$ \leerfeld & \tz $g(1)$ \leerfeld \\
\tz $f(-2)$ \leerfeld & \tz $g(-2)$ \leerfeld \\
\tz $h(-4)$ \leerfeld & \tz $f(1{,}5)$ \leerfeld \\
\tz $g(-0{,}5)$ \leerfeld & \\
\end{teilezwei}
\end{aufgabe}

\begin{aufgabe}{Ich kann prüfen, ob ein Punkt auf dem Graphen einer Potenzfunktion liegt.}
Prüfe mit der Punktprobe. Kreuze an.
\punktprobenliste{x^3/A(2|8), x^4/B(2|8), x^3/C({-3}|{-27}), x^4/D({-2}|{-16}), 0{,}5x^3/E(2|4), -x^4/F({-1}|1)}
\end{aufgabe}

\begin{aufgabe}{Ich kann zu einem Funktionswert das passende $x$ bestimmen.}
Bestimme alle $x$, die die Gleichung erfüllen.
\begin{teilezwei}
\tz $x^3 = 8$ \quad \feld{x} & \tz $x^3 = 1000$ \quad \feld{x} \\
\tz $x^3 = 1$ \quad \feld{x} & \tz $x^3 = 125$ \quad \feld{x} \\
\tz $x^4 = 16$ \quad \feld{x} & \tz $x^3 = -27$ \quad \feld{x} \\
\tz $x^4 = -1$ \quad \feld{x} & \\
\anweisung{Rechne mit dem Taschenrechner und runde auf zwei Nachkommastellen.}
\tz $x^3 = 50$ \quad \feld{x} & \\
\end{teilezwei}
\end{aufgabe}

\begin{aufgabe}{Ich kann am Funktionsterm erkennen, wie der Graph verläuft.}
Kreuze an, wie der Graph symmetrisch ist.
\begin{teile}
\teil $y = x^4$ \quad \kreuz{symmetrisch zur $y$-Achse}\kreuz{punktsymmetrisch zum Ursprung}
\teil $y = x^3$ \quad \kreuz{symmetrisch zur $y$-Achse}\kreuz{punktsymmetrisch zum Ursprung}
\teil $y = x^2$ \quad \kreuz{symmetrisch zur $y$-Achse}\kreuz{punktsymmetrisch zum Ursprung}
\teil $y = -2x^3$ \quad \kreuz{symmetrisch zur $y$-Achse}\kreuz{punktsymmetrisch zum Ursprung}
\end{teile}
\anweisung{Kreuze den Graphen an, der schmaler ist.}
\begin{teile}
\teil \kreuz{$y = 3x^4$}\quad\kreuz{$y = 0{,}5x^4$}
\end{teile}
\anweisung{Kreuze den Graphen an, der ganz unterhalb der $x$-Achse liegt (bis auf den Ursprung).}
\begin{teile}
\teil \kreuz{$y = x^4$}\quad\kreuz{$y = -x^4$}\quad\kreuz{$y = -x^3$}
\end{teile}
\end{aufgabe}

\begin{aufgabe}{Ich kann Gleichung und Graph einander zuordnen.}
Schreibe zu jeder Gleichung den Buchstaben des passenden Graphen.

\begin{ksys}[xmin=-2,xmax=2,ymin=-4,ymax=4,xstep=0.5,ablesen]
\funktion{\x^3}{}
\funktion{\x^4}{}
\funktion{-0.5*\x^4}{}
\funktion{0.25*\x^3}{}
\node[fill=white,inner sep=1pt] at (1.8,3.2) {$A$};
\node[fill=white,inner sep=1pt] at (-1.7,3.2) {$B$};
\node[fill=white,inner sep=1pt] at (1.85,-3.3) {$C$};
\node[fill=white,inner sep=1pt] at (1.75,2.2) {$D$};
\end{ksys}

\begin{teilezwei}
\tz $y = x^4$ \quad \leerfeld & \tz $y = x^3$ \quad \leerfeld \\
\tz $y = 0{,}25x^3$ \quad \leerfeld & \tz $y = -0{,}5x^4$ \quad \leerfeld \\
\end{teilezwei}
\end{aufgabe}

\begin{aufgabe}{Ich kann eine Potenzfunktion von einer Exponentialfunktion unterscheiden.}
Kreuze an. Rechne nichts.
\begin{teilezwei}
\tz $y = x^3$ \quad \kreuz{Potenz}\kreuz{Exponential} & \tz $y = 3^x$ \quad \kreuz{Potenz}\kreuz{Exponential} \\
\tz $y = 2x^4$ \quad \kreuz{Potenz}\kreuz{Exponential} & \tz $y = 4 \cdot 2^x$ \quad \kreuz{Potenz}\kreuz{Exponential} \\
\tz $y = x^2$ \quad \kreuz{Potenz}\kreuz{Exponential} & \tz $y = 0{,}5^x$ \quad \kreuz{Potenz}\kreuz{Exponential} \\
\end{teilezwei}
\end{aufgabe}

\begin{aufgabe}{Ich finde den Fehler bei Funktionswerten einer Potenzfunktion.}
Finde den Fehler, benenne ihn und korrigiere die Rechnung.
\begin{teile}
\teil Es ist $f(x) = -x^4$. Tim berechnet $f(2)$:
\rechnung{f(2) &= (-2)^4 \\ f(2) &= 16}
\teil Es ist $g(x) = x^3$. Sara berechnet $g(5)$:
\rechnung{g(5) &= 3 \cdot 5 \\ g(5) &= 15}
\end{teile}
\end{aufgabe}

\begin{aufgabe}{Ich kann begründen, wie der Graph einer Potenzfunktion verläuft.}
Begründe.
\begin{teile}
\teil Warum wird $y = x^4$ nie negativ?
\teil Warum verläuft der Graph von $y = x^3$ links unten durch den dritten Quadranten?
\end{teile}
\end{aufgabe}

\begin{aufgabe}{Ich kann mit einer Potenzfunktion eine Sachaufgabe zum Würfel lösen.}
Ein Würfel mit der Kantenlänge $k$ (in cm) hat das Volumen $V(k) = k^3$ (in cm³).
\begin{teile}
\teil Berechne das Volumen eines Würfels mit 5\,cm Kantenlänge.
\teil Ein würfelförmiges Gefäß fasst genau einen Liter (1000\,cm³). Wie lang ist seine Kante?
\teil Ein würfelförmiger Behälter soll 500\,cm³ fassen. Wie lang muss die Kante sein? Runde auf Millimeter.
\teil Wie ändert sich das Volumen, wenn man die Kantenlänge verdoppelt? Begründe.
\end{teile}
\end{aufgabe}

% abhaken.tex – erzeugt von gen_abhaken.py aus den Aufgabenquelltexten
\begin{abhakseite}
\ifmitzone
\abhakgruppe{Kennst du schon}
\abhak{1}{Ich kann Dezimalzahlen multiplizieren und das Ergebnis runden.}
\abhak{2}{Ich kann Punkte in ein Koordinatensystem eintragen.}
\abhak{3}{Ich kann Punkte eintragen, wenn eine Achse in größeren Schritten eingeteilt ist.}
\abhak{4}{Ich kann eine passende Einteilung für eine Achse wählen.}
\abhak{5}{Ich kann einen Wert um einen Prozentsatz erhöhen oder senken.}
\abhak{6}{Ich kann eine Zahlenfolge fortsetzen, wenn immer gleich viel dazukommt.}
\abhak{7}{Ich kann zwei Zahlen durcheinander teilen und das Ergebnis als Dezimalzahl angeben.}
\abhak{8}{Ich kann Prozentwert und Prozentsatz berechnen.}
\abhak{9}{Ich kann ein Guthaben mit Zinseszins berechnen.}
\abhak{10}{Ich kann Potenzen mit dem Taschenrechner berechnen.}
\abhak{11}{Ich finde den Fehler in einer Zinseszinsrechnung.}
\abhak{12}{Ich kann das Guthaben nach mehreren Jahren Zinseszins berechnen.}
\fi
\abhakgruppe{Einheit 1 von 5 · Lineares und exponentielles Wachstum unterscheiden}
\abhak{13}{Ich erkenne, ob jedes Mal derselbe Betrag oder derselbe Prozentsatz dazukommt.}
\abhak{14}{Ich erkenne, ob in einer Tabelle die Differenz oder der Quotient gleich bleibt.}
\abhak{15}{Ich erkenne, wo ein Graph die senkrechte Achse trifft.}
\abhakgruppe{\quad\textit{Wachstumsart an Tabelle und Text erkennen}}
\abhak{16}{Ich kann prüfen, ob in einer Tabelle immer derselbe Betrag dazukommt.}
\abhak{17}{Ich kann prüfen, ob in einer Tabelle immer mit derselben Zahl multipliziert wird.}
\abhak{18}{Ich kann an einer Tabelle entscheiden, ob ein Wert linear oder exponentiell wächst.}
\abhak{19}{Ich kann aus einem Text entscheiden, ob ein Wert linear oder exponentiell zu- oder abnimmt.}
\abhakgruppe{\quad\textit{Graph zu einem Wachstum auswählen}}
\abhak{20}{Ich erkenne, ob ein Graph eine Gerade oder eine gekrümmte Kurve ist.}
\abhak{21}{Ich kann den passenden Graphen zu einem Wachstum auswählen.}
\abhak{22}{Ich kann begründen, warum ein Graph nicht zu einem Wachstum passt.}
\abhakgruppe{\quad\textit{Wertepaare als Punkte darstellen}}
\abhak{23}{Ich kann Wertepaare eines Wachstums als Punkte eintragen.}
\abhak{24}{Ich kann Wertepaare eines Wachstums als Punkte eintragen – weiter.}
\abhak{25}{Ich kann für Wertepaare eine passende Achseneinteilung wählen.}
\abhak{26}{Ich finde den Fehler bei der Entscheidung zwischen linear und exponentiell.}
\abhak{27}{Ich kann begründen, warum bei gleichem Prozentsatz der Zuwachs immer größer wird.}
\abhak{28}{Ich kann Grenzen eines Wachstumsmodells nennen.}
\abhak{29}{Ich kann zwei Sparangebote mit linearem und exponentiellem Wachstum vergleichen.}
\abhakgruppe{Einheit 2 von 5 · Wachstumsfaktor und Wachstumstabelle}
\abhak{30}{Ich erkenne am Faktor, ob ein Wert zunimmt oder abnimmt.}
\abhakgruppe{\quad\textit{Wachstumsfaktor bestimmen}}
\abhak{31}{Ich kann aus einem Prozentsatz den Wachstumsfaktor bilden.}
\abhak{32}{Ich kann aus einem Wachstumsfaktor die prozentuale Änderung ablesen.}
\abhak{33}{Ich kann den Wachstumsfaktor als Quotient aus einer Tabelle bestimmen.}
\abhakgruppe{\quad\textit{Wachstumstabelle fortschreiben}}
\abhak{34}{Ich kann eine Wachstumstabelle mit dem Faktor fortschreiben.}
\abhak{35}{Ich kann eine Wachstumstabelle mit dem Faktor fortschreiben – weiter.}
\abhak{36}{Ich kann eine Wachstumstabelle mit dem Faktor fortschreiben – weiter.}
\abhak{37}{Ich finde den Fehler in einer Wachstumstabelle.}
\abhak{38}{Ich kann begründen, woran man den Wachstumsfaktor erkennt.}
\abhak{39}{Ich kann mit einer Wachstumstabelle eine Frage im Sachzusammenhang beantworten.}
\abhakgruppe{Einheit 3 von 5 · Exponentialfunktion aufstellen und auswerten}
\abhak{40}{Ich erkenne, welches Jahr der Schritt null ist.}
\abhakgruppe{\quad\textit{Wachstumsgleichung aufstellen und deuten}}
\abhak{41}{Ich kann eine Wachstumsgleichung aufstellen.}
\abhak{42}{Ich kann die passende Gleichung unter mehreren auswählen.}
\abhak{43}{Ich kann die Teile einer Wachstumsgleichung im Sachzusammenhang deuten.}
\abhakgruppe{\quad\textit{Werte und Schwellenwerte berechnen}}
\abhak{44}{Ich kann mit einer Wachstumsgleichung den Wert nach einigen Schritten berechnen.}
\abhak{45}{Ich kann eine Behauptung über ein Wachstum mit einer Rechnung prüfen.}
\abhak{46}{Ich kann berechnen, wann ein Wert eine Schwelle überschreitet.}
\abhak{47}{Ich finde den Fehler in einer Rechnung mit der Wachstumsgleichung.}
\abhak{48}{Ich kann begründen, warum in der Wachstumsgleichung eine Potenz steht.}
\abhak{49}{Ich kann mit Wachstumsgleichungen zwei Entwicklungen vergleichen.}
\abhakgruppe{Einheit 4 von 5 · Verdopplungs- und Halbwertszeit}
\abhak{50}{Ich erkenne, ob sich die Zeit oder der Bestand verdoppeln oder halbieren soll.}
\abhak{51}{Ich kann den Zielwert für die Verdopplung oder Halbierung angeben.}
\abhak{52}{Ich kann die Verdopplungs- oder Halbwertszeit aus einer Tabelle ablesen.}
\abhak{53}{Ich kann die Verdopplungs- oder Halbwertszeit am Graphen ablesen.}
\abhak{54}{Ich kann die Verdopplungs- oder Halbwertszeit durch Probieren mit dem Faktor bestimmen.}
\abhak{55}{Ich finde den Fehler bei der Verdopplungs- oder Halbwertszeit.}
\abhak{56}{Ich kann begründen, wann es eine feste Verdopplungszeit gibt.}
\abhak{57}{Ich kann mit der Halbwertszeit eine Frage im Sachzusammenhang beantworten.}
\abhakgruppe{Einheit 5 von 5 · Potenzfunktionen mit natürlichem Exponenten}
\abhakgruppe{\quad\textit{Wertetabelle und Graph}}
\abhak{58}{Ich kann eine Wertetabelle zu einer Potenzfunktion ausfüllen.}
\abhak{59}{Ich kann eine Wertetabelle zu einer Potenzfunktion ausfüllen – weiter.}
\abhak{60}{Ich kann den Graphen einer Potenzfunktion zeichnen.}
\abhakgruppe{\quad\textit{Funktionswerte, Punktprobe und Umkehrung}}
\abhak{61}{Ich kann Funktionswerte einer Potenzfunktion berechnen.}
\abhak{62}{Ich kann prüfen, ob ein Punkt auf dem Graphen einer Potenzfunktion liegt.}
\abhak{63}{Ich kann zu einem Funktionswert das passende $x$ bestimmen.}
\abhak{64}{Ich kann am Funktionsterm erkennen, wie der Graph verläuft.}
\abhak{65}{Ich kann Gleichung und Graph einander zuordnen.}
\abhak{66}{Ich kann eine Potenzfunktion von einer Exponentialfunktion unterscheiden.}
\abhak{67}{Ich finde den Fehler bei Funktionswerten einer Potenzfunktion.}
\abhak{68}{Ich kann begründen, wie der Graph einer Potenzfunktion verläuft.}
\abhak{69}{Ich kann mit einer Potenzfunktion eine Sachaufgabe zum Würfel lösen.}
\end{abhakseite}

\end{document}
